\documentclass[../../../include/open-logic-section]{subfiles}

\begin{document}

\olfileid{sth}{story}{blv}

\olsection{ضميمه: د فرېګه پنځم بنسټيز قانون}

په \olref[rus]{sec} کښې مو بيان کړل چې رسل خپل تناقض د هغه
نظام د ستونزې په توګه وړاندې کړے و چې فرېګه په خپل
\emph{Grundgesetze} کښې بيان کړے و۔ د فرېګه په نظام کښې د
ساده ادراک نېغ بيان نۀ و شامل۔ نو په دې ضميمه کښې به ډېر
لنډ بيان کړو چې د فرېګه نظام څه \emph{لرل}، له ساده ادراک
سره څنګه تړاو لري او د رسل له تناقض سره يې څه اړيکه ده۔

د فرېګه نظام \emph{د دويمې درجې} دے او د \emph{مفهوم د
غځونې} د تصور د بيانولو لپاره طرح شوے و۔\footnote{په کره
  توګه، فرېګه د لا عمومي تصور د رسمي کولو هڅه کوي: د تابع
  «د قيمتونو غځونه»۔ د مفهومونو غځونې د دې لا عمومي تصور
  ځانګړې بېلګه ده۔ د تفصيل لپاره \citet[pp.\ 8--9]{Heck2012}
  وګورئ۔} د فرېګه له نښو الهام اخيستے، $\fregeext{x}{F(x)}$
به د \emph{مفهوم~$F$ د غځونې} لپاره وليکو۔ دا يوه وسيله
ده چې يو \emph{پريديکات}، «$F$»، اخلي او په يوه (د لومړۍ
درجې) \emph{ترم}، «$\fregeext{x}{F(x)}$»، يې بدلوي۔ فرېګه
په دې وسيله د غړيتوب دا \emph{تعريف} وړاندې کړ:
\[
  a \in b =_\text{df} \exists G(b = \fregeext{x}{G(x)} \land Ga)
\]
په لنډه توګه: $a \in b$ که او يوازې که $a$ د داسې مفهوم
لاندې راځي چې غځونه يې $b$ ده۔ (ياد ولرئ چې کمي کوونکے
«$\exists G$» د دويمې درجې دے۔) فرېګه دا اصل هم مانه چې
د \emph{پنځم بنسټيز قانون} په نوم پېژندل کېږي:
$$\fregeext{x}{F(x)} = \fregeext{x}{G(x)} \liff \forall x (Fx \liff Gx)$$
په لنډه توګه: مفهومونه عين غځونې لري که او يوازې که پر عين
شيانو صدق کوي۔ (بيا هم، دواړه «$F$» او «$G$» د پريديکات
په ځاے کښې دي۔) اوس يو ساده اصل غړيتوب د خاصيت له
پوره کولو سره تړي:

\begin{lem}[په \emph{Grundgesetze} کښې]\ollabel{lem:Fregeextensions}
$\forall F \forall a(a \in \fregeext{x}{F(x)} \liff Fa)$
\end{lem}

\begin{proof} 
$F$ او $a$ ثابت ونيسئ۔ اوس $a \in \fregeext{x}{F(x)}$ که او يوازې که
$\exists G(\fregeext{x}{F(x)}
= \fregeext{x}{G(x)} \land Ga)$ (د غړيتوب د تعريف له مخې)، که او يوازې که
$\exists G(\forall x(Fx \liff Gx) \land Ga)$ (د پنځم بنسټيز قانون له مخې)، که او يوازې که $Fa$
(د دويمې درجې د بنسټيز منطق له مخې)۔
\end{proof}

او دا نژدې سملاسي ساده ادراک راکوي:

\begin{lem}[په \emph{Grundgesetze} کښې۔]
$\forall F \exists s \forall a (a \in s \liff Fa)$
\end{lem}

\begin{proof}
$F$ ثابت ونيسئ؛ اوس \olref{lem:Fregeextensions} دا راکوي چې $\forall a (a \in
\fregeext{x}{F(x)} \liff Fa)$؛ نو د وجودي عمومي کولو له مخې $\exists s\forall a(a \in s \liff Fa)$۔
نتيجه ترې راځي، ځکه $F$ هر يو ټاکل کېدلے شو۔
\end{proof}

که $F$ د $\forall x(Fx \liff x \notin x)$ له مخې ونيسو، د رسل تناقض ترې راځي۔

\end{document}
